% Figure: the ideal Ericsson cycle on the T-s plane. Two isotherms (at T_h and
% T_c) closed by two isobars, with the constant-pressure legs regenerated so the
% only external heat crosses the isotherms. Contrasts with the Stirling cycle
% (two isotherms closed by two isochores): both reach Carnot efficiency through
% an internal regenerator, differing in the shape of the connecting legs.
\begin{figure}[htbp]
\centering
\begin{tikzpicture}
\begin{axis}[
  width=11cm, height=7.5cm,
  xlabel={specific entropy $s$},
  ylabel={temperature $T$},
  xmin=0, xmax=6, ymin=200, ymax=1000,
  xtick=\empty, ytick={300,900},
  yticklabels={$T_c$,$T_h$},
  axis lines=left,
  tick label style={font=\small}, label style={font=\small},
  clip=false,
]
% Upper isotherm at T_h = 900 (heat in), from s=1.5 to s=4.5
\addplot[engred, very thick] coordinates {(1.5,900) (4.5,900)};
% Lower isotherm at T_c = 300 (heat out), from s=1.5 to s=4.5
\addplot[engblue, very thick] coordinates {(1.5,300) (4.5,300)};
% Right isobar: constant pressure cooling from (4.5,900) to (4.5,300):
% on T-s an isobar is a curve; approximate with a gentle sag to the left
\addplot[enggray, thick, samples=40, domain=300:900]
  ({4.5 - 0.55*ln(x/300)/ln(3)},{x});
% Left isobar: constant pressure heating from (1.5,300) to (1.5,900),
% regenerated (the mirror curve, offset so the two isobars are congruent)
\addplot[enggray, thick, samples=40, domain=300:900]
  ({1.5 - 0.55*ln(x/300)/ln(3)},{x});
% State labels
\node[font=\small, anchor=south] at (axis cs:1.5,900) {1};
\node[font=\small, anchor=south] at (axis cs:4.5,900) {2};
\node[font=\small, anchor=north] at (axis cs:4.5,300) {3};
\node[font=\small, anchor=north] at (axis cs:1.5,300) {4};
% Process annotations
\node[engred, font=\scriptsize, anchor=south] at (axis cs:3.0,900) {isothermal, $q_h$ in};
\node[engblue, font=\scriptsize, anchor=north] at (axis cs:3.0,300) {isothermal, $q_c$ out};
\node[enggray, font=\scriptsize, anchor=west, align=left] at (axis cs:2.7,600)
  {regenerated\\isobars};
\end{axis}
\end{tikzpicture}
\caption[Ideal Ericsson cycle on the temperature-entropy plane]{The ideal
Ericsson cycle: isothermal heat intake at $T_h$ (1 to 2) and isothermal heat
rejection at $T_c$ (3 to 4), closed by two constant-pressure legs (2 to 3 and
4 to 1). Because the two isobars are congruent, the heat shed on the cooling leg
exactly matches the heat wanted on the heating leg, so a counter-flow
regenerator can transfer it internally and the only heat crossing the cycle
boundary is the two isothermal exchanges. The efficiency is then $1 - T_c/T_h$,
the Carnot value, reached here through congruent isobars rather than the
Stirling cycle's congruent isochores. A real machine falls short by the
regenerator's incompleteness, the same term that governs the Stirling engine.}
\label{fig:vol04ch03:ericsson-ts}
\end{figure}
